Chương này giới thiệu bối cảnh nghiên cứu, động lực thực hiện báo cáo, phạm vi và mục tiêu, cũng như cấu trúc tổng thể của báo cáo.

\subsection{Bối cảnh và động lực}

Mạng di động thế hệ thứ năm (5G) không chỉ giới hạn ở cơ sở hạ tầng mặt đất. 3GPP trong Release 17 (2022) đã chính thức đưa hỗ trợ mạng phi mặt đất (Non-Terrestrial Network -- NTN) vào tiêu chuẩn New Radio (NR), cho phép các vệ tinh LEO, MEO và GEO hoạt động như các nút truy cập trong kiến trúc 5G~\cite{tr38821}. Xu hướng này mở ra khả năng phủ sóng toàn cầu -- đặc biệt cho các khu vực hải đảo, vùng sâu vùng xa và các ứng dụng hàng hải, hàng không -- vốn không thể được phục vụ bởi mạng mặt đất thông thường.

Tuy nhiên, kênh vệ tinh có hai đặc điểm cơ bản khác biệt so với kênh mặt đất:
\begin{itemize}
    \item \textbf{Độ trễ lan truyền lớn:} RTT từ 12{,}88 ms (LEO 600 km) trở lên theo TR 38.811~\cite{tr38811}, gấp hàng chục đến hàng trăm lần so với mạng mặt đất (thường $< 1$ ms).
    \item \textbf{Kênh LOS mạnh (Rician):} Do quỹ đạo vệ tinh ở độ cao lớn, thành phần tầm nhìn thẳng (LOS) chiếm ưu thế, được đặc trưng bởi hệ số Rician $K$ cao (điển hình 15 dB theo~\cite{tuninato2025}).
\end{itemize}

Cơ chế HARQ -- nền tảng của lớp MAC trong 5G NR -- được thiết kế cho RTT mặt đất và giới hạn tối đa 32 tiến trình (TS 38.214~\cite{ts38214}). Khi áp dụng cho NTN LEO, câu hỏi cốt lõi là: \textit{IR-HARQ còn hiệu quả ở mức độ nào, và cấu hình SCS nào giúp tránh pipeline stall?}

\subsection{Mục tiêu báo cáo}

Báo cáo này có ba mục tiêu chính:
\begin{enumerate}
    \item Xây dựng mô hình mô phỏng Monte Carlo cho IR-HARQ (Incremental Redundancy) trong kênh Rician NTN tại quỹ đạo LEO, tuân theo TR 38.811 và tham số hóa theo~\cite{tuninato2025}.
    \item Đánh giá định lượng hiệu quả IR-HARQ theo ba chiều: BLER vs SNR, ràng buộc N$_{\min}$ pipeline, và hiệu quả năng lượng trên bit.
    \item Xác định cấu hình SCS và số tiến trình HARQ phù hợp cho từng quỹ đạo LEO.
\end{enumerate}

\subsection{Phạm vi nghiên cứu}

\begin{itemize}
    \item \textbf{Kênh:} Rician phẳng một trạng thái (flat fading, single-state LMS), K = 15 dB, SCS 30 kHz (mặc định), $f_c = 20$ GHz, tốc độ UE 50 km/h.
    \item \textbf{Quỹ đạo:} LEO 600 km và LEO 1200 km (payload tái sinh, góc ngẩng 10°).
    \item \textbf{MCS:} MCS A (QPSK, $r = 1/2$).
    \item \textbf{HARQ:} IR-HARQ, tối đa 4 lần phát ($\text{max\_retx} = 3$), chuỗi RV \{0, 2, 3, 1\} (TS 38.212~\cite{ts38212}).
    \item \textbf{Không xét:} MIMO, beam management, overhead lớp cao.
\end{itemize}

\subsection{Cấu trúc báo cáo}

\begin{itemize}
    \item \textbf{Chương 2} trình bày cơ sở lý thuyết: tổng quan 5G-NTN, từ ARQ cơ bản đến IR-HARQ, đặc tính kênh vệ tinh Rician, và vấn đề pipeline trong NTN.
    \item \textbf{Chương 3} mô tả phương pháp mô phỏng: mô hình kênh, mô hình BLER dựa trên MI, bộ quản lý tiến trình HARQ và các tham số lựa chọn.
    \item \textbf{Chương 4} trình bày và phân tích kết quả mô phỏng từ ba thí nghiệm, bao gồm bảng số liệu và đồ thị.
    \item \textbf{Phần Kết luận} tóm tắt các phát hiện chính và đề xuất hướng phát triển.
\end{itemize}
